\documentclass[conference]{IEEEtran}
\IEEEoverridecommandlockouts
\usepackage{cite}
\usepackage{amsmath,amssymb,amsfonts}
\usepackage{algorithmic}
\usepackage{graphicx}
\usepackage{textcomp}
\usepackage{booktabs}
\usepackage{xcolor}
\usepackage{hyperref}
\usepackage{float}

\def\BibTeX{{\rm B\kern-.05em{\sc i\kern-.025em b}\kern-.08em
    T\kern-.1667em\lower.7ex\hbox{E}\kern-.125emX}}
\begin{document}


\title{Near-Real-Time ISAC-Based Jammer Detection and Localization using srsRAN\\}

\author{\IEEEauthorblockN{Samanvitha Balijepalli}
\IEEEauthorblockA{\textit{Dept. of ECE} \\
\textit{Rice University}\\
Houston, USA \\
sb273@rice.edu}
\and
\IEEEauthorblockN{Aniket Jadhav}
\IEEEauthorblockA{\textit{Dept. of ECE} \\
\textit{Rice University}\\
Houston, USA \\
aj156@rice.edu}
\and
\IEEEauthorblockN{Sandra Mansour}
\IEEEauthorblockA{\textit{Dept. of ECE} \\
\textit{Rice University}\\
Houston, USA  \\
sm367@rice.edu}
}

\maketitle

\begin{abstract}
Malicious jammers threaten wireless communications yet detecting them
typically requires costly dedicated RF monitoring hardware. This work
presents a software-based 5G framework for jammer detection and future
localization using the srsRAN stack without hardware modification. The
system reserves quiet symbol intervals within the 5G NR frame, creating
standards-compliant sensing windows at only $\approx$1.4\% scheduling
overhead. SNR and jammer power extracted from these windows feed
lightweight threshold-based classifiers. In a controlled ZMQ software
testbed, jammer power achieves perfect classification (F1\,=\,1.0) while
SNR-based detection reaches 92.18\% accuracy (F1\,=\,0.957). A
MUSIC-based angle-of-arrival baseline is separately validated, confirming
readiness for Phase~2 localization. Although results reflect a simulated
environment, they demonstrate feasible low-overhead jammer sensing and
establish a concrete foundation for near-real-time localization on USRP
hardware.
\end{abstract}

\begin{IEEEkeywords}
Integrated Sensing and Communication (ISAC), jammer detection, jammer
localization, srsRAN, 5G NR, wireless sensing, threshold-based detection,
MUSIC, angle-of-arrival
\end{IEEEkeywords}


% -----------------------------------------------------------------------
\section{Introduction}
% -----------------------------------------------------------------------


Malicious radio-frequency jammers pose an increasingly significant threat
to modern wireless systems. By flooding target bands with continuous-wave
or broadband interference, they can disable GPS receivers, degrade LTE/5G
base-station links, and disrupt safety-critical services such as emergency
communications and unmanned aerial vehicle control~\cite{3gpp38211}.
Conventional jammer detection relies on dedicated spectrum-monitoring
hardware — receivers, direction-finding antennas, and spectrum analysers —
that must be separately purchased, maintained, and re-tuned whenever
operating conditions change. This cost barrier motivates the development
of software-defined alternatives that reuse existing communication
infrastructure for sensing.

Integrated Sensing and Communication (ISAC) has emerged as a key
direction for next-generation wireless systems, enabling communication
infrastructure to simultaneously support sensing, detection, and
localization tasks~\cite{1143830,article}. By sharing radio resources
between data transmission and environmental awareness, ISAC improves
spectrum efficiency while eliminating the need for separate sensing
hardware. Open and programmable 5G software stacks, such as
srsRAN~\cite{srsran_docs} and Open5GS~\cite{open5gs_site}, provide a
particularly attractive experimental foundation because they expose
scheduler internals that commercial equipment conceals.


\subsection*{Related Work}

High-resolution subspace methods have long been the foundation of
direction-of-arrival (DoA) estimation. MUSIC (Multiple Signal
Classification), introduced by Schmidt~\cite{1143830}, decomposes the
array covariance matrix into signal and noise subspaces, yielding a
pseudospectrum with sharp peaks at true source angles. Root-MUSIC~\cite{45540}
and ESPRIT~\cite{32276} extend this framework by replacing grid search with
polynomial rooting or subspace rotation, reducing computation while
preserving resolution. Sparse reconstruction methods~\cite{inbook} offer
further gains under limited snapshots but at higher algorithmic complexity.
SpotFi~\cite{article} demonstrated that OFDM channel state information
(CSI) — already collected during normal Wi-Fi communication — can feed a
modified MUSIC pipeline for indoor localization, achieving sub-meter
accuracy without additional hardware. A key limitation of SpotFi is that
it assumes cooperative access points and a static, well-characterized
indoor environment. Our work targets the more adversarial setting of
jammer localization in a 5G NR environment, where the interference
source is uncooperative, the signal structure is unknown, and
software-defined gNB modifications are required to create sensing
opportunities in the first place.

This work presents a software-based framework for jammer detection and
future localization using srsRAN within an open 5G communication
environment. Rather than relying on dedicated sensing signals or external
RF instrumentation, the proposed system creates standards-compliant sensing
opportunities directly from the communication process by reserving quiet
symbol intervals within the 5G NR frame. As illustrated in
Fig.~\ref{fig:system_framework}, the current phase focuses on jammer
detection using communication-quality indicators such as SNR and jammer
power, while the future phase extends the framework toward localization
using richer signal features and the MUSIC algorithm.

The main contributions of this work are: (i) a low-overhead,
scheduler-based sensing design that creates quiet windows inside the 5G
NR frame with only $\approx$1.4\% resource overhead; (ii) an evaluation
of threshold-based jammer detectors showing that jammer power achieves
perfect separation in the current dataset; and (iii) a validated
MUSIC-based angle-of-arrival baseline that serves as the localization
reference for the next phase. Together these results establish a practical
foundation for near-real-time jammer localization in an open-source 5G
platform.

\begin{figure}[H]
    \centering
    \includegraphics[width=0.48\textwidth]{Abstract.png}
    \caption{Joint communication and sensing framework for jammer detection
    and localization using data-scheduling-based sensing opportunities in a
    5G gNB. Phase~1 (current work) performs detection; Phase~2+ will
    estimate jammer position using frequency-domain and time-of-flight
    features fed into the MUSIC estimator.}
    \label{fig:system_framework}
\end{figure}


% -----------------------------------------------------------------------
\section{Methodology}
% -----------------------------------------------------------------------



\subsection{System Architecture}

The proposed system is built on the srsRAN~\cite{srsran_docs} 5G software
stack and uses communication-driven sensing opportunities to detect and
eventually localize jammer activity without dedicated sensing hardware.
The full platform consists of an Open5GS~\cite{open5gs_site} core network,
a modified srsRAN gNB, an srsUE instance, and a jammer proxy inserted
into the ZeroMQ (ZMQ) signal path, as shown in Fig.~\ref{fig:system_flowchart}.

The gNB scheduler was modified to blank two OFDM symbols per sensing
period, creating periodic \emph{quiet windows} during which the transmitter
is silent and the receiver captures ambient RF energy, interference, and
jammer-induced signal components. IQ samples from these windows are stored
for post-processing. Simultaneously, the jammer proxy logs UE power,
jammer power, and signal-to-interference ratio (SIR) values synchronized
to the same slot timeline, providing the ground-truth labels needed for
detection evaluation.

A post-processing layer converts the captured IQ data and signal logs into
MATLAB-compatible data structures. In the current stage, sensing is
performed using communication-quality indicators derived from the scheduled
windows — particularly SNR and jammer power — processed by lightweight
threshold-based detectors. The longer-term goal is to extend this pipeline
toward localization using richer signal features and the MUSIC
algorithm~\cite{1143830}, consistent with prior structured-measurement
approaches such as SpotFi~\cite{article}.

\begin{figure}[htbp]
    \centering
    \includegraphics[width=0.8\linewidth]{System Architecture.png}
    \caption{System architecture and processing flow for the proposed
    sensing framework.}
    \label{fig:system_flowchart}
\end{figure}

\subsection{Data-Scheduling-Based Sensing Design}

Rather than relying on Sounding Reference Signal (SRS) sweeps, the
framework introduces sensing opportunities directly through data
scheduling within the 5G NR frame, motivated by 3GPP TS~38.211
Table~4.3.2-1~\cite{3gpp38211}, which designates symbols 12 and 13 as
flexible symbols in standard TDD slot formats. When configured as downlink
quiet symbols, these intervals let the gNB observe the channel without
expecting UE transmission.

A new configuration structure was added to srsRAN to expose
sensing-related parameters: enable flag, quiet symbol start index, quiet
symbol count, and sensing period in slots. At 30~kHz subcarrier spacing
one slot lasts 0.5~ms, and reserving 2 symbols once every 10 slots yields
a sensing overhead of
\[
\text{Overhead} = \frac{2}{10 \times 14} = \frac{2}{140} \approx 1.4\%.
\]
This small resource cost creates periodic sensing opportunities within the
normal communication process.

During each quiet sensing window, IQ samples are captured from the receive
path and written to disk. The jammer proxy simultaneously injects
controlled continuous-wave (CW) interference and logs UE power, jammer
power, and SIR. These synchronized measurements are the input to the
detection stage described next.

\subsection{Detection and Localization Processing Framework}

A MATLAB post-processing framework was developed to analyze the sensing
measurements captured during the scheduled quiet windows. The main
variables in the current stage are time, measured SNR, jammer power, and
a ground-truth jammer-state label (OFF / ON).

To quantify class separability, summary statistics were computed. The
sample mean and standard deviation are
\begin{equation}
\mu = \frac{1}{N}\sum_{i=1}^{N} x_i, \qquad
\sigma = \sqrt{\frac{1}{N-1}\sum_{i=1}^{N}(x_i-\mu)^2}.
\end{equation}

Using class means, the midpoint thresholds are
\begin{equation}
T_{\mathrm{SNR}} = \frac{\mu_{\mathrm{SNR,OFF}} + \mu_{\mathrm{SNR,ON}}}{2}, \quad
T_{\mathrm{JP}} = \frac{\mu_{\mathrm{JP,OFF}} + \mu_{\mathrm{JP,ON}}}{2}.
\end{equation}

The three detectors are defined as
\begin{align}
\hat{y}_{\mathrm{SNR}}(i) &=
\begin{cases} 1, & \mathrm{SNR}(i) < T_{\mathrm{SNR}} \\ 0, & \text{otherwise,} \end{cases}
\label{eq:snr_det} \\
\hat{y}_{\mathrm{JP}}(i) &=
\begin{cases} 1, & P_J(i) > T_{\mathrm{JP}} \\ 0, & \text{otherwise,} \end{cases}
\label{eq:jp_det} \\
\hat{y}_{\mathrm{C}}(i) &=
\begin{cases} 1, & \hat{y}_{\mathrm{SNR}}(i)=1 \;\lor\; \hat{y}_{\mathrm{JP}}(i)=1 \\ 0, & \text{otherwise.} \end{cases}
\label{eq:combined_det}
\end{align}

Overall accuracy is
\begin{equation}
\mathrm{Acc} = \frac{TP + TN}{TP + TN + FP + FN}.
\label{eq:accuracy}
\end{equation}


Because the dataset is class-imbalanced (501 OFF vs.\ 8953 ON samples,
roughly a 1:18 ratio), accuracy alone can be misleading. We therefore
also report precision $P = TP/(TP+FP)$, recall $R = TP/(TP+FN)$, and
F1-score $F_1 = 2PR/(P+R)$ for each detector.

To support the future localization stage, a separate framework evaluates
MUSIC-based angle-of-arrival (AoA) estimation using a 3-element uniform
linear array (ULA). Spatial covariance matrix decomposition separates the
signal and noise subspaces, after which the MUSIC pseudospectrum is
evaluated for angle estimation. This stage is not yet integrated into the
srsRAN sensing path but provides a validated reference model consistent
with prior high-resolution methods~\cite{1143830,45540,32276,inbook}.

\subsection{srsRAN Test Environment}

The full system was deployed on a Linux server using the srsRAN 5G
software stack, Open5GS core network, and ZMQ-based simulated radio links.
The system was configured for NR Band~3 (center frequency 1842.5~MHz,
20~MHz bandwidth, 30~kHz subcarrier spacing, single-antenna TDD,
single-user scenario).

The gNB and UE communicated normally during standard operation. The jammer
proxy introduced a configurable CW interference signal after slot~2000,
injecting at 100~kHz frequency offset with amplitude~0.8. Runtime
execution was monitored via command-line logs to verify gNB
initialization, UE attachment, and jammer activation.


\textit{Testbed limitation:} ZMQ-based transport introduces padding values
into the receive buffer during certain quiet-symbol intervals. Some
captured IQ samples therefore reflect simulation transport behavior rather
than true ambient RF. This does not invalidate the threshold-based
detection results — which rely on the jammer proxy's synchronized log
values rather than raw IQ power — but it does mean the present results
should be interpreted as a controlled software testbed study. A future
USRP hardware implementation would capture actual ambient RF and jammer
energy directly.


% -----------------------------------------------------------------------
\section{Results}
% -----------------------------------------------------------------------

\subsection{Jammer Injection and Measurement Overview}

Experiments were conducted to evaluate jammer detection using the proposed
scheduler-based sensing pipeline. A total of 9454 sensing measurements
were recorded: 501 jammer OFF samples (slots 0--1999) and 8953 jammer ON
samples (slots 2000 onward), as summarized in Table~\ref{tab:dataset_stats}.

Before jammer activation the average measured SNR was approximately
$11.07$~dB. After activation the SNR dropped sharply to approximately
$-91.10$~dB, indicating severe communication degradation. Simultaneously,
measured jammer power increased from 0 to approximately 0.64. These
feature shifts establish a clear distinction between the two states and
motivate the threshold-based detection results presented below.

\begin{table}[htbp]
\centering
\caption{Dataset statistics under jammer OFF and jammer ON conditions.}
\label{tab:dataset_stats}
\begin{tabular}{l r}
\hline
\textbf{Metric} & \textbf{Value} \\
\hline
Total samples & 9454 \\
Jammer OFF samples & 501 \\
Jammer ON samples & 8953 \\
Class ratio (OFF:ON) & 1:18 \\
SNR mean (OFF) & 11.07 dB \\
SNR std.\ (OFF) & 31.23 dB \\
SNR mean (ON) & $-$91.10 dB \\
SNR std.\ (ON) & 26.34 dB \\
Jammer power mean (OFF) & 0.00 \\
Jammer power std.\ (OFF) & 0.00 \\
Jammer power mean (ON) & 0.64 \\
Jammer power std.\ (ON) & 0.00 \\
\hline
\end{tabular}
\end{table}

\subsection{Feature Behavior Over Time}

Fig.~\ref{fig:jammer_power_time} shows jammer power over time, split by
jammer state. The OFF samples cluster near zero; the ON samples cluster
near 0.64 with no visible overlap. Fig.~\ref{fig:snr_time} shows the
measured SNR over time. During normal operation SNR is high and positive;
after jammer activation it collapses to strongly negative values, with
intermittent spikes caused by repeated UE reconnection attempts.

\begin{figure}[htbp]
    \centering
    \includegraphics[width=0.8\linewidth]{Jammer Power on:off.png}
    \caption{Jammer power versus time for jammer OFF and jammer ON conditions.}
    \label{fig:jammer_power_time}
\end{figure}

\begin{figure}[htbp]
    \centering
    \includegraphics[width=0.8\linewidth]{SNR_time.png}
    \caption{SNR versus time under jammer sensing conditions.}
    \label{fig:snr_time}
\end{figure}

\subsection{Threshold-Based Jammer Detection}

Three threshold-based detectors were evaluated; numerical performance is
summarized in Table~\ref{tab:detection_performance}.


\begin{table}[htbp]
\centering
\caption{Performance of threshold-based jammer detectors.}
\label{tab:detection_performance}
\resizebox{\columnwidth}{!}{%
\begin{tabular}{l r r r r r r r r r}
\hline
\textbf{Detector} & \textbf{Thresh.} & \textbf{Acc.} & \textbf{Prec.} & \textbf{Rec.} & \textbf{F1} & \textbf{TP} & \textbf{TN} & \textbf{FP} & \textbf{FN} \\
\hline
SNR-based         & $-$40.02 dB & 92.18\% & 100\% & 91.74\% & 0.957 & 8214 & 501 & 0 & 739 \\
Jammer-power      & 0.320       & \textbf{100\%}   & 100\% & 100\%   & 1.000 & 8953 & 501 & 0 & 0 \\
Combined          & N/A         & \textbf{100\%}   & 100\% & 100\%   & 1.000 & 8953 & 501 & 0 & 0 \\
\hline
\end{tabular}%
}
\end{table}

Fig.~\ref{fig:snr_detection} shows the SNR-based detector. The jammer
OFF and ON SNR distributions are substantially separated, but some spread
near the threshold leaves the classes imperfectly separable. The detector
achieved 92.18\% accuracy with 739 false negatives and an F1-score of
0.957. The false negatives correspond to intermittent SNR spikes during
UE reconnection attempts, which temporarily push jammer-ON measurements
above the threshold.

Fig.~\ref{fig:jammer_power_detection} shows the jammer-power-based
detector. The OFF and ON distributions are completely disjoint: all OFF
values sit at exactly 0 and all ON values sit at 0.64. The detector
achieved 100\% accuracy with F1 = 1.0 and no errors, providing the
strongest first-stage detection feature.

The combined detector (Eq.~\eqref{eq:combined_det}) also achieves 100\%
accuracy and F1 = 1.0, but does not improve on the jammer-power detector
alone. This is expected: jammer power already provides perfect separation
in the current dataset, so adding the SNR criterion contributes nothing
extra. Its value would likely emerge under more challenging conditions
(see Section~\ref{sec:discussion}).

\begin{figure}[htbp]
    \centering
    \includegraphics[width=0.8\linewidth]{SNR Jammer Detection.png}
    \caption{SNR-based jammer detection histogram. The vertical red line
    marks the decision threshold at $-$40.02~dB.}
    \label{fig:snr_detection}
\end{figure}

\begin{figure}[htbp]
    \centering
    \includegraphics[width=0.8\linewidth]{Jammer Power Detection.png}
    \caption{Jammer-power-based detection histogram. The two classes are
    completely separated with no overlap.}
    \label{fig:jammer_power_detection}
\end{figure}

\subsection{MUSIC Localization Baseline}

To support the future localization stage, a separate MATLAB simulation
validated a MUSIC-based AoA baseline using a 3-element ULA.
Fig.~\ref{fig:music_result} presents four diagnostics:
\begin{itemize}
    \item \textbf{CSI phase:} Approximately stable across subcarriers for
          all three antennas, indicating a well-conditioned input with
          preserved spatial phase structure.
    \item \textbf{I-Q constellation:} Tight clustering confirms acceptable
          baseband signal integrity.
    \item \textbf{Covariance eigenvalues:} One dominant eigenvalue,
          consistent with a rank-1 single-source assumption.
    \item \textbf{MUSIC pseudospectrum:} Clear dominant peak aligned with
          the $25^\circ$ ground-truth angle.
\end{itemize}
Together these results confirm that the MUSIC estimator correctly isolates
the signal subspace~\cite{1143830} and provides a valid localization
baseline for Phase~2 integration.

\begin{figure}[htbp]
    \centering
    \includegraphics[width=0.95\linewidth]{MUSIC_Simulation_result.png}
    \caption{MUSIC-based angle-of-arrival estimation baseline: CSI phase
    (top-left), I-Q constellation (bottom-left), covariance eigenvalues
    (top-right), and MUSIC pseudospectrum with peak at $25^\circ$
    (bottom-right).}
    \label{fig:music_result}
\end{figure}


% -----------------------------------------------------------------------
\section{Discussion}
\label{sec:discussion}
% -----------------------------------------------------------------------

\textbf{Detection performance.}
Jammer activity produces strong, consistent feature changes: mean SNR
drops from 11.07~dB to $-$91.10~dB and mean jammer power rises from 0 to
0.64 (Table~\ref{tab:dataset_stats}). The jammer-power detector achieves
perfect classification (F1 = 1.0), while the SNR detector achieves F1 =
0.957. The 739 SNR false negatives arise from UE reconnection spikes that
temporarily elevate measured SNR during the jammer-ON period, reflecting
the greater variability of SNR as a communication-quality metric.

\textbf{Class imbalance.}
The dataset contains roughly 18$\times$ more jammer-ON than jammer-OFF
samples (Table~\ref{tab:dataset_stats}), a direct consequence of the
experimental design: jammer activation was delayed to slot~2000 to
establish a clean baseline, leaving a short OFF window. Because the
jammer-power detector has zero FP and zero FN, the imbalance does not
affect its evaluation; all metrics equal 1.0 regardless of class weights.
The SNR detector is similarly unaffected by imbalance in the FP direction
(zero FP), and precision remains 100\%. Future experiments should balance
the class distribution more carefully or apply stratified cross-validation
to ensure that results generalize across different ON/OFF ratios.



\textbf{Testbed limitations.}
As noted in Section~II-D, the ZMQ transport introduces IQ padding
artifacts during quiet-symbol intervals. The present detection results
are not invalidated by this — they rely on the proxy log values rather
than raw IQ power — but a hardware deployment on USRP radios is required
before the IQ-domain features needed for MUSIC-based localization can be
captured reliably from actual RF.


\textbf{Interpreting the 100\% accuracy result.}
Perfect classification in a detection experiment warrants scrutiny. In
the current setup, jammer power is logged directly by the proxy — it
reflects the injected signal amplitude rather than an estimated RF
quantity derived from IQ data. The separation is therefore exact by
construction: OFF samples carry exactly 0 power and ON samples carry
exactly 0.64. This result demonstrates that the \emph{measurement
pipeline} reliably captures jammer state, but it does not yet establish
robustness against realistic impairments such as multipath fading,
antenna pattern variation, background interference, or jammers that
modulate their power over time. SNR, with its 92.18\% accuracy, is a
better proxy for true RF-domain difficulty.

\textbf{Combined detector.}
The combined detector matches jammer-power performance exactly,
confirming that jammer power dominates in the current experiment. This
does not imply SNR is irrelevant: in a hardware deployment with fading or
when the jammer amplitude varies, SNR may provide complementary evidence
that improves robustness. A weighted fusion or machine-learning classifier
could exploit both features more flexibly in future work.

\textbf{MUSIC baseline and Phase 2 roadmap.}
The stable CSI phase, tight I-Q clustering, dominant covariance
eigenvalue, and 25$^\circ$-aligned pseudospectrum peak confirm that the
MUSIC estimator is correctly isolating the signal subspace~\cite{1143830}.
The planned Phase~2 integration proceeds in three steps: (i) port IQ
capture from ZMQ to USRP hardware; (ii) extract per-subcarrier CSI
matrices from the quiet-window samples; and (iii) feed a multi-antenna
MUSIC pipeline with those CSI matrices to estimate jammer AoA in near
real time. Combining AoA from two or more gNBs will then yield a
2-D jammer position estimate.

\textbf{Summary.}
Communication-driven sensing opportunities created through scheduling are
sufficient for reliable first-stage jammer detection in an open-source 5G
framework. Jammer power is the most discriminative feature under the
present conditions, while SNR remains a useful secondary feature whose
value will likely grow in more challenging RF environments. These findings
validate the low-overhead sensing architecture and establish a concrete
roadmap for extending the system to near-real-time jammer localization.

The entire source code and related images can be found at \href{https://github.com/AniketttJadhav/srsRAN-Open5GS}{this GitHub repository}.

% -----------------------------------------------------------------------
\section*{Acknowledgments}
% -----------------------------------------------------------------------

This work was conducted as part of the Master of Electrical and Computer
Engineering program at Rice University. The authors acknowledge the
academic support and institutional resources provided by Rice University,
as well as the guidance and feedback from Dr.\ Rahman Doost throughout
the project. The statements made in this work are solely the
responsibility of the authors.



\section*{CRediT Author Statement}

\textbf{Sandra Mansour:} Software, Formal analysis, Visualization,
Writing---Original draft. \textbf{Samanvitha Balijepalli:}
Investigation, Algorithm analysis, Writing---Review \& Editing.
\textbf{Aniket Jadhav:} Investigation, Methodology, Software.
\textbf{Dr.\ Rahman Doost:} Conceptualization, Supervision, Project
administration, Resources, Writing---Review \& Editing.

\bibliographystyle{IEEEtran}
\bibliography{references}

\end{document}